% 机械波（高中）
% license Usr
% type Tutor

\subsection{机械波}
机械波是机械振动在空间中传播的现象。称振动的源头为波源，称承载振动的物质为介质\footnote{真空实际并非空无一物，而是能量最低的一个状态。实际上，从量子物理的观点来看，无论真空能是否为0（我们可以灵活地调整能量零点），真空总是在极短时间内产生和湮灭粒子-反粒子对。这种量子涨落引起声波传播已被实验验证}。从微观上看，只要有波源在作振动，且质点之间有弹性恢复力，那么必然引起传播方向上，下一个相邻质点作振动。每一个质点，都可以看作接下来的，机械振动的波源，如果没有其他能量损耗，机械波的传播是没有止境的。其传播路径上的每个质点都可视为弹簧振子，在受到上一个质点的扰动后，才开始在平衡位置附近作振动

从上述过程来看，机械波对振动能量的传播不言而喻。实际上，机械波还是振动状态（相位）的传播。如果A点为较近点，B点是较远点，且振动从A传播到B需要时间$t$，那么B在某时刻$t_0$的相位总是B在$t_0-t$的相位。
\subsubsection{波的图像}
如果我们把机械波简化成一连串质点的运动，且设波源作简谐振动，介质是理想介质（即波速稳定），那么在任意时刻，给这一串质点“拍照”，它们的位置构成简谐函数的波形。
比如设波源的振动形式为$y=\sin(\omega t)$，波速为$v$，那么距波源$x$的质点位移为
\begin{equation}
y=\sin \omega (t-\frac{x}{v})~.
\end{equation}
